%=====================================================================
\chapter{Amortised Analysis}\label{ch:amortised}
%=====================================================================
\locator{4}{Part IV --- The Design Paradigms}

\begin{outcomes}
By the end of this chapter you will be able to:
\begin{tight}
  \item \textbf{Explain} why multiplying a worst-case operation cost by the
        number of operations can give a badly loose bound.
  \item \textbf{Apply} the aggregate, accounting and potential methods to a
        sequence of operations \emph{(CLO-3)}.
  \item \textbf{Prove} that $n$ appends to a doubling array cost $\bigTh{n}$ in
        total.
  \item \textbf{Distinguish} amortised cost from average-case and worst-case
        cost.
  \item \textbf{Judge} when an amortised guarantee is adequate and when it is
        not.
\end{tight}
\end{outcomes}

\begin{prereq}
\chref{cases} for the three cases --- amortised cost is a fourth thing and must
be kept apart from all of them. \chref{structures} for the \algname{Build-Heap}
analysis, which is the same looseness in miniature.
\end{prereq}

\begin{keyterms}
amortised analysis $\bullet$ aggregate method $\bullet$ accounting method
$\bullet$ credit $\bullet$ potential method $\bullet$ potential function
$\bullet$ dynamic array $\bullet$ growth factor $\bullet$ table doubling
$\bullet$ worst case versus amortised
\end{keyterms}

\begin{alertbox}[Why this chapter is here, when HEC does not name it]
The published outline does not list amortised analysis. It is taught anyway,
for three reasons, and the Course Specification in the front matter records
them:

\begin{tightnum}
  \item \chref{hashing}'s claim that a hash table keeps $\alpha$ bounded
        depends on resizing, whose cost is amortised. Without this chapter that
        claim is an assertion.
  \item \chref{mst}'s union--find is the standard example of a structure whose
        cost is \emph{only} meaningful amortised.
  \item CLO-3 asks for the time complexity of simple algorithms, and for a
        sequence of operations on a growing structure that is an amortised
        question. Answering it per-operation gives the wrong answer.
\end{tightnum}
\end{alertbox}

\section{A Loose Bound}

\begin{conceptbox}[The problem, stated]
Appending to a dynamic array is almost always $\bigO{1}$: write at the end,
increment the length. Occasionally the buffer is full, and then the append
costs $\bigTh{n}$ --- allocate a bigger buffer and copy everything.

\smallskip
So the worst case of a single append is $\bigTh{n}$, and
\[ n \text{ appends} \;\le\; n \times \bigTh{n} = \bigTh{n^{2}}. \]

\smallskip
That bound is \textbf{correct and useless}. It is correct because no append
exceeds $\bigO{n}$; it is useless because the expensive appends cannot all
happen --- after one, the buffer is twice as big and cannot fill again for
another $n$ appends.

\smallskip
\term{Amortised analysis} bounds the \emph{total} cost of a sequence directly,
without pretending every operation is the worst one.
\end{conceptbox}

\begin{definitionbox}[Amortised Cost]
The \term{amortised cost} of an operation is the total cost of a worst-case
sequence of $n$ operations, divided by $n$.

\smallskip
It is a \textbf{worst-case} quantity, not an average over a distribution. No
probability is involved and no assumption about the input is made --- which is
what distinguishes it sharply from the average case of \chref{cases} and from
the expected time of \chref{quicksort}.
\end{definitionbox}

\begin{alertbox}[Four costs, kept apart]
\rowcolors{2}{white}{lightbg}
\begin{longtable}{@{}L{2.8cm}L{5.4cm}L{5.8cm}@{}}
\toprule
\rowcolor{primary!15}
\textbf{Cost} & \textbf{Averaged over} & \textbf{Assumes} \\
\midrule
\endhead
worst case & nothing --- it is a maximum & nothing \\
average case & a distribution on inputs & that distribution \\
expected & the algorithm's own coin flips & nothing about the input \\
\textbf{amortised} & a \emph{sequence} of operations & nothing at all \\
\bottomrule
\end{longtable}

\smallskip
An amortised bound is as strong as a worst-case one: it holds for every
sequence, with no probability anywhere. What it does \emph{not} promise is that
any individual operation is fast, and Section~17.4 measures exactly how slow
one of them gets.
\end{alertbox}

\section{Three Methods, One Answer}

\dsfig{\aoaalg{\algname{Append}$(A, x)$}{
\If{$A.\mathit{len} = A.\mathit{cap}$}
  \State $\mathit{cap}' \gets \max(1, 2 \cdot A.\mathit{cap})$
  \State allocate $B$ of size $\mathit{cap}'$
  \State copy $A[0 \mathinner{\ldotp\ldotp} A.\mathit{len}]$ into $B$
    \Comment{the expensive step}
  \State $A.\mathit{buf} \gets B$;\ \ $A.\mathit{cap} \gets \mathit{cap}'$
\EndIf
\State $A[A.\mathit{len}] \gets x$;\ \ $A.\mathit{len} \gets A.\mathit{len}+1$
}}{\algname{Append} with table doubling. Line 4 runs rarely, and when it does
it costs $\bigTh{n}$. The three methods below are three ways of saying that the
rarity exactly pays for the cost.}{append}

\subsection*{The aggregate method: add it all up}

\begin{theorem}[Doubling array]\label{thm:doubling}
Starting from an empty array, $n$ calls to \algname{Append} cost $\bigTh{n}$ in
total, so the amortised cost per append is $\bigTh{1}$.
\end{theorem}

\begin{proof}
Copying happens only when the length reaches the capacity, which with doubling
occurs at lengths $1, 2, 4, 8, \ldots$ --- the powers of two up to $n$. The
copy at length $2^{k}$ moves $2^{k}$ elements. So the total copying over $n$
appends is
\[ \sum_{k=0}^{\lfloor \lg n \rfloor} 2^{k} \;=\; 2^{\lfloor\lg n\rfloor+1}-1
   \;<\; 2n. \]
Adding the $n$ constant-time writes gives total cost under $3n = \bigTh{n}$,
and dividing by $n$ gives $\bigTh{1}$ amortised.
\end{proof}

\begin{conceptbox}[The accounting method: pay in advance]
Charge each append \textbf{3 units} of amortised cost, and spend them like
this:

\begin{tight}
  \item 1 unit pays for writing the new element --- the actual work;
  \item 1 unit is saved as \term{credit} on the new element;
  \item 1 unit is saved as credit on an element in the older half of the array.
\end{tight}

\smallskip
When the array of capacity $m$ fills, it last doubled at size $m/2$, so the
$m/2$ most recent elements have each saved 2 units: $m$ units of credit in all,
exactly enough to pay for copying $m$ elements.

\smallskip
The credit never goes negative, so the total charged ($3n$) is an upper bound
on the total spent, giving $\bigO{1}$ amortised. This is the same accounting a
business does when it sets aside money each month against an annual bill.
\end{conceptbox}

\begin{conceptbox}[The potential method: a bank balance]
Define a \term{potential function} $\Phi$ on the data structure's state ---
roughly, how much trouble it is currently storing up. For the dynamic array,
\[ \Phi = 2\,\mathit{len} - \mathit{cap}. \]
The amortised cost of an operation is
\[ \hat{c}_{i} = c_{i} + \Phi_{i} - \Phi_{i-1}, \]
its real cost plus the change in potential.

\smallskip
\textbf{Cheap append} (no resize): $c_{i} = 1$, $\mathit{len}$ rises by 1 so
$\Phi$ rises by 2, giving $\hat{c}_{i} = 1 + 2 = 3$.

\smallskip
\textbf{Expensive append} (resize from $m$ to $2m$, with $\mathit{len} = m$
before): $c_{i} = m + 1$. Before, $\Phi = 2m - m = m$; after,
$\Phi = 2(m+1) - 2m = 2$. So
\[ \hat{c}_{i} = (m+1) + 2 - m = 3. \]

\smallskip
\textbf{Both cases give exactly 3.} The potential was built up by the cheap
operations and spent by the expensive one. Since $\Phi \ge 0$ always (the array
is at least half full after any doubling) and $\Phi_{0} = 0$, the total real
cost is at most the total amortised cost, $3n$.

\smallskip
\textbf{Why bother with a third method?} Because the potential method is the
one that generalises. The aggregate method needs you to see the whole sequence;
the accounting method needs a story about who pays. The potential method needs
only a function, and it is how the union--find bound in \chref{mst} and the
Fibonacci heap bound in \chref{shortestpaths} are proved.
\end{conceptbox}

\section{Measured}

\begin{worked}[The Total, Counted Exactly]
\texttt{code/ch17\_amortised.cpp} counts every element copy.

\rowcolors{2}{white}{lightbg}
\begin{center}
\begin{tabular}{@{}rrrrr@{}}
\toprule
\rowcolor{primary!15}
$n$ & doubling copies & $/n$ & grow-by-one copies & $/(n^{2}/2)$ \\
\midrule
 1{,}000 &  1{,}023 & 1.023 &        499{,}500 & 0.9990 \\
 8{,}000 &  8{,}191 & 1.024 &     31{,}996{,}000 & 0.9999 \\
32{,}000 & 32{,}767 & 1.024 &    511{,}984{,}000 & 1.0000 \\
64{,}000 & 65{,}535 & 1.024 &  2{,}047{,}968{,}000 & 1.0000 \\
\bottomrule
\end{tabular}
\end{center}

\smallskip
\textbf{Doubling: copies$/n$ is 1.024, flat.} Theorem~\ref{thm:doubling}
predicted under 2; the actual constant is just over 1, because
$\sum_{k<\lg n} 2^{k} = n-1$ when $n$ is a power of two.

\smallskip
\textbf{Grow-by-one: exactly $n^{2}/2$.} The quotient is $0.9990 \to 1.0000$,
confirming $\sum_{k=1}^{n-1} k = n(n-1)/2$ to four decimal places.

\smallskip
\textbf{So the $\bigTh{n^{2}}$ bound was not paranoia} --- it is exactly right
for the wrong growth policy. What amortised analysis shows is that it is
\emph{not} right for doubling, and the difference between the two policies is
one line of code.
\end{worked}

\begin{worked}[Amortised $\bigO{1}$ Does Not Mean Every Operation Is Fast]
One million appends, doubling:

\rowcolors{2}{white}{lightbg}
\begin{center}
\begin{tabular}{@{}lr@{}}
\toprule
\rowcolor{primary!15}
quantity & value \\
\midrule
appends & 1{,}000{,}000 \\
reallocations & 21 \\
appends that copied anything & 20 \quad(0.0020\%) \\
total copies & 1{,}048{,}575 \quad(1.049 per append) \\
\textbf{most expensive single append} & \textbf{524{,}288 copies} \\
worst case $\times\, n$ (the useless bound) & 524{,}288{,}000{,}000 \\
\bottomrule
\end{tabular}
\end{center}

\smallskip
\textbf{Read the last two rows together.} The amortised cost is 1.049 copies
per append. The worst single append copied 524{,}288 elements --- half a
million times the amortised figure.

\smallskip
\textbf{Both numbers are true and they describe different things.} If you are
computing throughput, 1.049 is the number you want. If you are meeting a
latency budget --- a frame deadline, an interrupt handler, a trading system ---
then 1.049 is irrelevant and 524{,}288 is the number that will end your
career.

\smallskip
\textbf{This is the chapter's main caution.} Amortised $\bigO{1}$ is a
guarantee about totals. When individual operations must be bounded, the answer
is a different structure --- an incrementally resized array that copies a few
elements on every append, trading a worse constant for a real worst-case
bound.

\smallskip
\textbf{And note how rare the expensive ones are:} 20 of a million appends,
0.002\%. A latency histogram with 1{,}000 samples would very likely contain
none of them, and the system would look fine until it was not.
\end{worked}

\begin{worked}[The Growth Factor Is a Real Decision]
$n = 1{,}000{,}000$ appends under different growth factors:

\rowcolors{2}{white}{lightbg}
\begin{center}
\begin{tabular}{@{}rrrr@{}}
\toprule
\rowcolor{primary!15}
factor & total copies & per append & wasted capacity \\
\midrule
1.10 & 10{,}170{,}704 & 10.171 &  1.7\% \\
1.25 &  4{,}489{,}047 &  4.489 & 12.2\% \\
1.50 &  2{,}099{,}753 &  2.100 &  5.0\% \\
2.00 &  1{,}048{,}575 &  1.049 &  4.9\% \\
3.00 &    797{,}161 &  0.797 & \textbf{59.4\%} \\
4.00 &    349{,}525 &  0.350 &  4.9\% \\
\bottomrule
\end{tabular}
\end{center}

\smallskip
\textbf{Every factor above 1 gives $\bigTh{1}$ amortised} --- the sum is
geometric with ratio $1/f$ and converges for any $f > 1$. The asymptotics
cannot distinguish these rows at all.

\smallskip
\textbf{What distinguishes them is the constant and the memory.} A factor of
1.1 does ten times as much copying as doubling; a factor of 3 wastes 59\% of
its memory at this particular $n$. (The waste figures are lumpy because they
depend on where $n$ falls relative to the growth sequence --- 2.0 and 4.0 both
land near a power of two here.)

\smallskip
\textbf{Several standard libraries use 1.5 rather than 2}, and the reason is
not in this table: with a factor of 2, the new block is always larger than the
sum of all previously freed blocks, so the allocator can never reuse them. With
1.5 it eventually can. That is an argument about the allocator, not about the
algorithm, and it is a good example of a decision that asymptotic analysis
cannot reach.
\end{worked}

\begin{worked}[And in Seconds]
\rowcolors{2}{white}{lightbg}
\begin{center}
\begin{tabular}{@{}rrr|rr@{}}
\toprule
\rowcolor{primary!15}
$n$ & doubling (ms) & ratio & $n$ & grow-by-one (ms) \\
\midrule
 1{,}000{,}000 &   3.5 & ---  &  20{,}000 &    84.1 \\
 4{,}000{,}000 &  14.7 & 2.02 &  40{,}000 &   350.5 \\
16{,}000{,}000 &  58.2 & 1.93 &  80{,}000 & 1{,}432.8 \\
32{,}000{,}000 & 119.3 & 2.05 & 160{,}000 & 5{,}759.3 \\
\bottomrule
\end{tabular}
\end{center}

\smallskip
Doubling's ratios average 2.03 --- linear. Grow-by-one's average 4.09 ---
quadratic. At $n = 160{,}000$ the grow-by-one version takes 5.8 seconds; the
doubling version handles 32 million elements in 119 milliseconds.

\smallskip
Note the two columns use different ranges of $n$. Doubling at $n = 160{,}000$
takes under a millisecond, which is too small to measure --- the first version
of this experiment reported ratios of 0.92, 1.19, 1.81 and 4.04 for a linear
algorithm, all of it noise.
\end{worked}

\section{Where Else This Pattern Appears}

\begin{conceptbox}[A list worth keeping]
\rowcolors{2}{white}{lightbg}
\begin{longtable}{@{}L{3.8cm}L{3.4cm}L{3.0cm}L{3.4cm}@{}}
\toprule
\rowcolor{primary!15}
\textbf{Structure} & \textbf{Operation} & \textbf{Worst case} &
\textbf{Amortised} \\
\midrule
\endhead
dynamic array & append & $\bigTh{n}$ & $\bigTh{1}$ \\
hash table with resizing & insert & $\bigTh{n}$ & $\bigTh{1}$ \\
binary counter & increment & $\bigTh{\lg n}$ & $\bigTh{1}$ \\
union--find & find & $\bigTh{\lg n}$ & $\bigO{\alpha(n)}$ \\
splay tree & search & $\bigTh{n}$ & $\bigTh{\lg n}$ \\
\algname{Build-Heap} & whole operation & --- & $\bigTh{n}$ not
  $\bigTh{n\lg n}$ \\
\bottomrule
\end{longtable}

\smallskip
The last row is \chref{structures}'s result, and it is the same phenomenon:
multiplying a count by a worst case gave $\bigO{n\lg n}$, and the truth was
$\bigTh{n}$ because the expensive cases are rare.

\smallskip
The union--find row is \chref{mst}'s, where $\alpha(n)$ is the inverse
Ackermann function --- below 5 for any $n$ that can be written down. It is the
most dramatic amortised bound in common use, and it is proved by the potential
method.
\end{conceptbox}

\begin{pitfall}
\begin{tight}
  \item \textbf{Multiplying worst case by count.} Measured, that bound was
        524{,}288{,}000{,}000 against a true cost of 1{,}048{,}575 --- a
        factor of half a million.
  \item \textbf{Calling amortised cost an average case.} There is no
        distribution and no probability; it is a worst-case bound on a total.
  \item \textbf{Promising low latency from an amortised bound.} The worst
        single append copied 524{,}288 elements.
  \item \textbf{Growing by a constant amount.} Measured exactly $n^{2}/2$
        copies. The policy, not the structure, is what makes it quadratic.
  \item \textbf{A potential function that can go negative.} The argument
        requires $\Phi \ge \Phi_{0}$, or the total is not bounded.
  \item \textbf{Assuming the growth factor does not matter because the
        asymptotics agree.} Measured, 1.1 did ten times the copying of 2.0.
  \item \textbf{Benchmarking with too few samples to see the rare case.} 0.002\%
        of appends were expensive; most histograms will miss them entirely.
\end{tight}
\end{pitfall}

\begin{lab}[The Rare Expensive Operation]
Reproduces all five tables. Expect twenty-five minutes.

\begin{lstlisting}[language=C++]
// ch17_amortised.cpp -- a sequence of operations, costed as a sequence.
// Build: cl /O2 /EHsc /std:c++17 ch17_amortised.cpp  (see Appendix A)

// --- 1. A dynamic array with an adjustable growth policy -----------
// Deliberately hand-rolled so every copy can be counted. "factor" is
// the multiplier on resize; a factor of 0 means "grow by one", which
// is the policy the quadratic analysis describes.
    void push(int x) {
        if (len == cap) {
            size_t ncap = (factor <= 1.0) ? cap + 1
                        : std::max<size_t>(1, (size_t)(cap * factor));
            if (ncap == cap) ncap = cap + 1;
            std::vector<int> nb(ncap);
            for (size_t i = 0; i < len; ++i) { nb[i] = buf[i]; ++copies; }
            buf.swap(nb);
            cap = ncap;
            ++reallocs;
        }
        buf[len++] = x;
    }

// --- 2. Total copies: the aggregate method, measured ---------------
// Doubling: the copies are 1 + 2 + 4 + ... + n/2 < n, so the total
// work for n appends is under 2n and the amortised cost per append is
// under 2. Grow-by-one: 1 + 2 + ... + (n-1) = n(n-1)/2.

// --- 3. Amortised O(1) does not mean every operation is cheap ------
// Print the cost of the single most expensive append, and how many
// appends are expensive at all. This is the distinction between
// amortised and worst-case that the chapter exists for.
        long long worst = *std::max_element(cost.begin(), cost.end());
        size_t expensive = 0;
        for (long long c : cost) if (c > 0) ++expensive;

// --- 4. The growth factor is a real engineering choice -------------
// Larger factor: fewer copies, more wasted space. Smaller factor: the
// reverse. 1.5 is what several standard libraries use; 2.0 is what
// most textbooks describe.

// --- 5. And in seconds, against the real thing ---------------------
// The two policies need different ranges. Doubling at n = 160,000
// takes under a millisecond, too small to measure -- the first
// version of this table reported ratios of 0.92, 1.19, 1.81 and 4.04
// for a linear algorithm, all of it noise.
\end{lstlisting}

\textbf{What to notice.} Part 3 prints \texttt{worst case x n} beside the true
total. Those two numbers --- 524 billion and 1.05 million --- are the bound
amortised analysis replaces and the answer it gives. The whole chapter is the
gap between them.

\smallskip
\textbf{Then try to break it.} Add a growth factor of exactly 1.0 and predict
what happens before running. The guard \texttt{if (ncap == cap) ncap = cap + 1}
catches it and the policy degenerates to grow-by-one --- but remove that guard
and the capacity never increases, the \texttt{push} never terminates its
resize, and you have written \chref{invariants}'s non-terminating loop again. A
measure that fails to decrease is still the commonest way to hang a program.
\end{lab}

\begin{checkpoint}
\begin{tightnum}
  \item Using the aggregate method, show that $n$ appends to a doubling array
        perform fewer than $2n$ copies, and compare with the measured 1.024 per
        append.
  \item Compute the amortised cost of an append under the potential function
        $\Phi = 2\,\mathit{len} - \mathit{cap}$, for both the cheap and the
        expensive case, and say why they must come out equal.
  \item Give a situation in which an amortised $\bigO{1}$ guarantee is
        inadequate, and say what the measured data in Section~17.3 tells you
        about it \emph{(CLO-3)}.
\end{tightnum}
\end{checkpoint}

\begin{chaptersummary}
\begin{tight}
  \item Multiplying a worst-case operation cost by the number of operations can
        be enormously loose, because the expensive operations cannot all
        happen.
  \item Amortised cost is the total cost of a worst-case \emph{sequence},
        divided by its length. It assumes nothing about inputs and involves no
        probability --- unlike average-case or expected cost.
  \item Three methods: aggregate (sum it), accounting (charge extra and save
        credit), potential (define $\Phi$ and let $\hat c = c + \Delta\Phi$).
        The potential method is the one that generalises.
  \item $n$ appends to a doubling array cost $\bigTh{n}$: the copies are
        $1+2+4+\cdots < 2n$. Measured exactly 1.024 copies per append.
  \item Measured: grow-by-one costs exactly $n^{2}/2$ copies. The policy makes
        the difference, not the structure.
  \item Measured: of one million appends, 20 copied anything, the amortised
        cost was 1.049, and the worst single append copied 524{,}288 elements.
        Amortised $\bigO{1}$ says nothing about any individual operation.
  \item Any growth factor $> 1$ gives $\bigTh{1}$ amortised; measured, a factor
        of 1.1 did ten times the copying of 2.0, and 3.0 wasted 59\% of its
        memory.
\end{tight}
\end{chaptersummary}

\begin{reviewq}
  \item \textbf{(MCQ)} The amortised cost of an append to a doubling array is:
        (a)~$\bigTh{1}$; (b)~$\bigTh{\lg n}$; (c)~$\bigTh{n}$;
        (d)~$\bigTh{n^{2}}$.
  \item \textbf{(MCQ)} Amortised analysis differs from average-case analysis in
        that it: (a)~assumes a uniform input distribution; (b)~assumes nothing
        about the input; (c)~applies only to randomised algorithms; (d)~gives a
        lower bound.
  \item \textbf{(Short)} State the three methods of amortised analysis and give
        the one-line idea of each \emph{(CLO-3)}.
  \item \textbf{(Short)} Explain why an amortised $\bigO{1}$ guarantee may be
        unacceptable for a real-time system, citing the measured worst single
        append.
  \item \textbf{(Short)} Why does any growth factor greater than 1 give
        $\bigTh{1}$ amortised appends, while growing by a constant does not?
  \item \textbf{(Applied)} Analyse the amortised cost of incrementing a binary
        counter $n$ times, using both the aggregate and the potential method
        (take $\Phi$ = the number of 1 bits).
  \item \textbf{(Applied)} A team reports their append-heavy service has a
        99.99th-percentile latency 500$\times$ its median. Explain the likely
        cause from this chapter, say what measurement would confirm it, and
        give two fixes with their trade-offs.
\end{reviewq}
